\section*{Bab \ref{ch:g4:geometry} --- Garis dan Segi Banyak}

\begin{solution}{exo:g4:geometry:1}
Lukisan mengikuti \cref{met:g4:geometry:draw}; persegi kecil
menandai \omterm{def:g3:shapes:rightangle}{sudut siku-siku} antara $d$ dan garis tegak lurusnya.
\end{solution}

\begin{solution}{exo:g4:geometry:2}
E: ketiga goresan mendatarnya \omterm{def:g4:geometry:def}{sejajar}; masing-masing \omterm{def:g4:geometry:def}{tegak lurus}
terhadap goresan tegaknya. H: kedua goresan tegaknya \omterm{def:g4:geometry:def}{sejajar}, palang
melintangnya \omterm{def:g4:geometry:def}{tegak lurus} terhadap keduanya. N: kedua goresan
tegaknya \omterm{def:g4:geometry:def}{sejajar}; goresan miringnya bukan keduanya. T: kedua
goresannya \omterm{def:g4:geometry:def}{tegak lurus}. Z: goresan atas dan bawahnya \omterm{def:g4:geometry:def}{sejajar};
goresan miringnya bukan keduanya.
\end{solution}

\begin{solution}{exo:g4:geometry:3}
\omterm{def:g4:geometry:polygon}{Segi enam}; \omterm{def:g4:geometry:polygon}{segi empat}; segitiga; \omterm{def:g4:geometry:polygon}{segi delapan}; \omterm{def:g4:geometry:polygon}{segi lima}. \omterm{def:g4:geometry:polygon}{Segi
banyak} punya \omterm{def:g4:geometry:polygon}{titik sudut} sebanyak sisinya: $6$, $4$, $3$, $8$, $5$.
\end{solution}

\begin{solution}{exo:g4:geometry:4}
\omterm{def:g4:geometry:polygon}{Segi lima} punya $5$ sisi. Dari $A$, yang bukan tetangganya adalah
$C$ dan $D$: dua diagonal, yaitu $[AC]$ dan $[AD]$.
\end{solution}

\begin{solution}{exo:g4:geometry:5}
$OA = 4$ cm (\omterm{def:g3:shapes:shapes}{jari-jari}); $OM = 4$ cm (setiap titik pada lingkaran);
$AB = 8$ cm (\omterm{def:g4:geometry:circle}{diameter} dua kali \omterm{def:g3:shapes:shapes}{jari-jari}).
\end{solution}

\begin{solution}{exo:g4:geometry:6}
Bukalah jangka sepanjang satu sisi, lalu bawalah bukaan itu ke dua
sisi yang lain: kalau rentang jarum dan pensilnya pas dengan sisi
kedua, kedua sisi itu sama panjang.
\end{solution}

\begin{solution}{exo:g4:geometry:7}
Jalur itu menutup menjadi \omterm{def:g4:geometry:polygon}{segi banyak} dengan $5$ sisi --- sebuah
\omterm{def:g4:geometry:polygon}{segi lima} (yang tidak beraturan!).
\end{solution}

\begin{solution}{exo:g4:geometry:8}
``Dua \omterm{def:g4:geometry:def}{garis tegak lurus} selalu berpotongan'': benar --- keduanya
berpotongan tepat di \omterm{def:g3:shapes:rightangle}{sudut siku-sikunya}.

``Dua \omterm{def:g4:geometry:def}{garis sejajar} tidak pernah berpotongan'': benar --- itulah definisinya.

``Sebuah garis dapat \omterm{def:g4:geometry:def}{sejajar} dengan satu garis dan \omterm{def:g4:geometry:def}{tegak lurus}
terhadap garis yang lain'': benar --- di kertas berpetak, garis
mendatar \omterm{def:g4:geometry:def}{sejajar} dengan setiap garis mendatar dan \omterm{def:g4:geometry:def}{tegak lurus} terhadap setiap garis tegak.
\end{solution}

\begin{solution}{exo:g4:geometry:9}
Garis $d$ dan $f$ tampak \omterm{def:g4:geometry:def}{sejajar} --- dan memang selalu \omterm{def:g4:geometry:def}{sejajar}: dua
garis yang \omterm{def:g4:geometry:def}{tegak lurus} terhadap garis yang sama pasti \omterm{def:g4:geometry:def}{sejajar}
(\cref{prop:g6:lines:facts} meresmikannya).
\end{solution}

\begin{solution}{exo:g4:geometry:10}
Satu program yang benar: ``1. Gambarlah persegi $ABCD$ dengan sisi
$4$ cm. 2. Tandailah titik $O$ tempat kedua diagonalnya berpotongan.
3. Gambarlah lingkaran berpusat $O$ dengan \omterm{def:g3:shapes:shapes}{jari-jari} $2$ cm.''
Lingkaran itu menyentuh setiap sisi di tengahnya.
\end{solution}

\begin{solution}{exo:g4:geometry:11}
\omterm{def:g4:geometry:polygon}{Segi empat}: $2$ diagonal. \omterm{def:g4:geometry:polygon}{Segi lima}: $5$. \omterm{def:g4:geometry:polygon}{Segi enam}: $9$. Setiap
\omterm{def:g4:geometry:polygon}{titik sudut} baru menambah diagonal makin banyak --- hitungannya
tumbuh makin cepat ($+3$, lalu $+4$, \dots).

\end{solution}
